\documentclass{article}
 
%%%%%%%%%%%%%%%%%%%%%%%%%%%%%%%%%%%%%%%%%%%%%%%%%%%%%%%%%%%%%%%%%%%%%%%%%%
%%%%%%%%%%%%%%%%%%%%%			USEPACKAGE			%%%%%%%%%%%%%%%%%%%%%
%%%%%%%%%%%%%%%%%%%%%%%%%%%%%%%%%%%%%%%%%%%%%%%%%%%%%%%%%%%%%%%%%%%%%%%%%%
\usepackage[top=3cm, bottom=3cm, outer=3cm, inner=3cm]{geometry}
\usepackage{multicol}
\usepackage{graphicx}
\usepackage{url}
\usepackage{xcolor}
\usepackage{hyperref}
\hypersetup{
	colorlinks=true,
	linkcolor=blue,
	filecolor=magenta,
	urlcolor=cyan,
	citecolor=green,
}
\usepackage{array}
\newcolumntype{x}[1]{>{\centering\arraybackslash\hspace{0pt}}p{#1}}
\usepackage{natbib}
\usepackage{pdfpages}
\usepackage{multirow}
\usepackage[normalem]{ulem}
\useunder{\uline}{\ul}{}
\usepackage{listings}
\usepackage{listingsutf8}
\usepackage{amsmath,amssymb}
\usepackage{color, colortbl}
\usepackage[english,spanish]{babel}
\usepackage[utf8]{inputenc}
\usepackage{algpseudocode}
\usepackage{algorithm}
\usepackage{verbatim}
\usepackage{caption}
\usepackage{subcaption}
\usepackage{float}
 
\AtBeginDocument{\selectlanguage{spanish}}
\renewcommand{\figurename}{Figura}
\renewcommand{\refname}{Referencias}
\AtBeginDocument{\renewcommand\tablename{Tabla}}
 
\usepackage[T1]{fontenc}
\usepackage{imakeidx}
\makeindex[columns=3, title=Índice Alfabético, intoc]
 
%%%%%%%%%%%%%%%%%%%%%%%%%%%%%%%%%%%%%%%%%%%%%%%%%%%%%%%%%%%%%%%%%%%%%%%%%%
%%%%%%%%%%%%%%%%%%%%%			DEFINE COLOR			%%%%%%%%%%%%%%%%%%%%%
%%%%%%%%%%%%%%%%%%%%%%%%%%%%%%%%%%%%%%%%%%%%%%%%%%%%%%%%%%%%%%%%%%%%%%%%%%
\definecolor{dkgreen}{rgb}{0,0.6,0}
\definecolor{gray}{rgb}{0.5,0.5,0.5}
\definecolor{mauve}{rgb}{0.58,0,0.82}
\definecolor{tablebackground}{rgb}{0.0, 0.45, 0.63}
 
\lstdefinestyle{ascii-tree}{
    literate={├}{|}1 {─}{--}1 {└}{+}1
}
 
\lstset{
	extendedchars=true,
	inputencoding=utf8,
	basicstyle={\small\ttfamily},
	showstringspaces=false,
	commentstyle=\color{dkgreen},
	keywordstyle=\color{blue},
	frame=single,
	framextopmargin=4pt,
	framexbottommargin=4pt,
	framexleftmargin=4pt,
	framexrightmargin=2pt,
	language=Java,
	aboveskip=3mm,
	belowskip=3mm,
	columns=flexible,
	numbers=none,
	numberstyle=\tiny,
	stringstyle=\color{mauve},
	breaklines=true,
	breakatwhitespace=true,
	tabsize=2,
	rulecolor=\color{black!30},
	backgroundcolor=\color{gray!10},
}
 
%%%%%%%%%%%%%%%%%%%%%%%%%%%%%%%%%%%%%%%%%%%%%%%%%%%%%%%%%%%%%%%%%%%%%%%%%%
%%%%%%%%%%%%%%%%%%%%%			NEW COMMAND			%%%%%%%%%%%%%%%%%%%%%
%%%%%%%%%%%%%%%%%%%%%%%%%%%%%%%%%%%%%%%%%%%%%%%%%%%%%%%%%%%%%%%%%%%%%%%%%%
\newcommand{\itemEmail}{cvalienteq@ulasalle.edu.pe}
\newcommand{\itemTeacher}{Carmen Aida Valiente Quiroz}
\newcommand{\itemAdvisor}{Richart Escobedo Quispe}
\newcommand{\itemCourse}{Lenguaje de Programación I}
\newcommand{\itemUniversity}{Universidad La Salle}
\newcommand{\itemFaculty}{Facultad de Ingenierías y Arquitectura}
\newcommand{\itemDepartment}{Departamento Académico de Ingeniería y Matemáticas}
\newcommand{\itemSchool}{Carrera Profesional de Ingeniería de Software}
\newcommand{\itemAcademic}{2026 - I}
\newcommand{\itemInput}{05 Julio 2026}
\newcommand{\itemOutput}{06 Julio 2026}
\newcommand{\itemPracticeNumber}{Final}
\newcommand{\itemTheme}{Tres en Raya Quiz}
 
\usepackage{fancyhdr}
\pagestyle{fancy}
\fancyhf{}
\setlength{\headheight}{30pt}
\renewcommand{\headrulewidth}{1pt}
\renewcommand{\footrulewidth}{1pt}
\fancyhead[L]{\raisebox{0.2\height}{\includegraphics[width=3.0cm]{img/logo_internacional_secundario_color.png}}}
\fancyhead[C]{\fontsize{7}{7}\selectfont \itemUniversity \\ \itemFaculty \\ \itemDepartment \\ \itemSchool \\ \textbf{\itemCourse}}
\fancyfoot[L]{\itemTeacher}
\fancyfoot[C]{Pág. \thepage}
\fancyfoot[R]{\itemCourse}
 
%%%%%%%%%%%%%%%%%%%%%%%%%%%%%%%%%%%%%%%%%%%%%%%%%%%%%%%%%%%%%%%%%%%%%%%%%%
%%%%%%%%%%%%%%%%%%%%%			BEGIN				%%%%%%%%%%%%%%%%%%%%%
%%%%%%%%%%%%%%%%%%%%%%%%%%%%%%%%%%%%%%%%%%%%%%%%%%%%%%%%%%%%%%%%%%%%%%%%%%
\begin{document}
 
\vspace*{10px}
 
\begin{center}
 \fontsize{17}{17} \textbf{ Examen Final (Parte I)}
\end{center}
\centerline{\textbf{\Large Tema: \itemTheme}}
 
%%% TABLE 01 HEAD
\begin{table}[H]
	\begin{tabular}{|x{4.7cm}|x{4.8cm}|x{4.8cm}|}
		\hline
		\rowcolor{tablebackground}
		\color{white} \textbf{Estudiante} & \color{white}\textbf{Carrera Profesional} & \color{white}\textbf{Asignatura} \\
		\hline
		{\itemTeacher \par \itemEmail} & \itemSchool & {\itemCourse} \\
		\hline
	\end{tabular}
\end{table}
 
%%% TABLE 02 HEAD
\begin{table}[H]
	\begin{tabular}{|x{3cm}|x{6.5cm}|x{4.8cm}|}
		\hline
		\rowcolor{tablebackground}
		\color{white}\textbf{Proyecto} & \color{white}\textbf{Asesor} & \color{white}\textbf{Duración} \\
		\hline
		\itemPracticeNumber & \itemAdvisor & 12 horas \\
		\hline
	\end{tabular}
\end{table}
 
%%% TABLE 03 HEAD
\begin{table}[H]
	\begin{tabular}{|x{4.7cm}|x{4.8cm}|x{4.8cm}|}
		\hline
		\rowcolor{tablebackground}
		\color{white}\textbf{Semestre académico} & \color{white}\textbf{Fecha de inicio} & \color{white}\textbf{Fecha de entrega} \\
		\hline
		\itemAcademic & \itemInput & \itemOutput \\
		\hline
	\end{tabular}
\end{table}
 
\tableofcontents
\clearpage
 
%%%%%%%%%%%%%%%%%%%%%%%%%%%%%%%%%%%%%%%%%%%%%%%%%%%%%%%%%%%%%%%%%%%%%%%%%%
%%%%%%%%%%%%%%%%%%%%%		ESPECIFICACIONES		%%%%%%%%%%%%%%%%%%%%%
%%%%%%%%%%%%%%%%%%%%%%%%%%%%%%%%%%%%%%%%%%%%%%%%%%%%%%%%%%%%%%%%%%%%%%%%%%
\section{Especificaciones del Proyecto}
 
\subsection{Objetivos del proyecto}
\begin{itemize}
	\item Analizar e implementar la arquitectura MVC básica en Java para un juego por turnos.
	\item Diseñar y construir una interfaz de usuario interactiva y moderna mediante la API de Java Swing.
	\item Integrar avatares dinámicos mediante gráficos vectoriales con colores personalizados en 2D.
	\item Implementar un sistema de trivia académica que promueva la retroalimentación pedagógica ante errores.
	\item Diseñar un criterio objetivo de desempate basado en el tiempo de respuesta.
\end{itemize}
 
\subsection{Equipos, materiales y temas}
\begin{itemize}
	\item Sistema Operativo {\tt Windows 11}: Versión 25H2
	\item Editor de texto o IDE {\tt Visual Studio Code} con extensiones Java.
	\item JDK (Java Development Kit) versión 17 o superior.
	\item Control de versiones con {\tt Git} y repositorio en {\tt GitHub}.
\end{itemize}
 
%%%%%%%%%%%%%%%%%%%%%%%%%%%%%%%%%%%%%%%%%%%%%%%%%%%%%%%%%%%%%%%%%%%%%%%%%%
%%%%%%%%%%%%%%%%%%%%%		MARCO TEÓRICO			%%%%%%%%%%%%%%%%%%%%%
%%%%%%%%%%%%%%%%%%%%%%%%%%%%%%%%%%%%%%%%%%%%%%%%%%%%%%%%%%%%%%%%%%%%%%%%%%
\section{Marco Teórico}
 
\subsection{Java Swing}
\begin{itemize}
	\item Java Swing es una biblioteca gráfica para Java que proporciona un conjunto de componentes GUI ligeros para crear interfaces independientes de la plataforma.
	\item Utiliza una arquitectura orientada a componentes donde se anidan paneles (\texttt{JPanel}), botones (\texttt{JButton}) y etiquetas (\texttt{JLabel}) dentro de contenedores principales como el \texttt{JFrame}.
	\item El renderizado vectorial se logra mediante \texttt{Graphics2D}, permitiendo dibujar figuras geométricas, trazados y gradientes de color suavizados por antialiasing.
\end{itemize}
 
\subsection{Programación Orientada a Objetos en Java}
\begin{itemize}
	\item La aplicación hace uso extensivo de conceptos clave de POO como herencia (extendiendo de \texttt{JFrame}), encapsulación (atributos de clase con modificadores de acceso private), y polimorfismo mediante el manejo de interfaces como \texttt{ActionListener}.
\end{itemize}
 
\subsection{Síntesis de Sonido MIDI}
\begin{itemize}
	\item En lugar de depender de archivos pesados y de formatos incompatibles (\texttt{.wav}), el proyecto utiliza la biblioteca nativa \texttt{javax.sound.midi}.
	\item Esto nos permite inicializar un sintetizador (\texttt{Synthesizer}) y programar notas individuales y acordes de manera asíncrona, asegurando portabilidad absoluta sin dependencias externas.
\end{itemize}

\subsection{Especificaciones de la interfaz del juego}
\begin{table}[H]
	\caption{Especificaciones de interfaz y lógica}
	\setlength{\tabcolsep}{0.5em}
	{\renewcommand{\arraystretch}{1.5}
	\begin{tabular}{|p{5cm}|p{2cm}|x{5cm}|x{2cm}|}
		\hline
		\textbf{Resolución ventana} & 680x790 & \textbf{Layout del tablero} & Grid 3x3 \\
		\hline
		\textbf{Sonido integrado} & MIDI nativo & \textbf{Carga de Avatares} & JFileChooser / Vector \\
		\hline
		\textbf{Temas de preguntas} & POO y Swing & & \\
		\hline
	\end{tabular}
	}
\end{table}

\begin{figure}[H]
	\centering
	\includegraphics[width=0.85\textwidth,keepaspectratio]{img/IMAGEN1.png}
	\caption{Captura del juego finalizado con visualización del tablero y tarjetas}
\end{figure}
 
\clearpage
 
%%%%%%%%%%%%%%%%%%%%%%%%%%%%%%%%%%%%%%%%%%%%%%%%%%%%%%%%%%%%%%%%%%%%%%%%%%
%%%%%%%%%%%%%%%%%%%%%		GUÍA DE TRABAJO			%%%%%%%%%%%%%%%%%%%%%
%%%%%%%%%%%%%%%%%%%%%%%%%%%%%%%%%%%%%%%%%%%%%%%%%%%%%%%%%%%%%%%%%%%%%%%%%%
\section{Estructura del Proyecto y Trabajo}
 
\subsection{Repositorio Local y Git}
Para inicializar el proyecto en la consola de comandos se utilizaron las siguientes instrucciones de Git:

\begin{lstlisting}[language=bash, caption={Clonación y estado del proyecto}]
$ git clone https://github.com/cvalienteq-coder/lp1.git
$ cd lp1/TresEnRayaGUI
$ git status
\end{lstlisting}

\subsection{Archivo .gitignore}
El archivo \texttt{.gitignore} se configuró para evitar subir archivos de clase binarios generados por el compilador de Java:
\begin{lstlisting}[language=bash, caption={Archivo .gitignore}]
*.class
.project
.classpath
\end{lstlisting}
 
%%%%%%%%%%%%%%%%%%%%%%%%%%%%%%%%%%%%%%%%%%%%%%%%%%%%%%%%%%%%%%%%%%%%%%%%%%
%%%%%%%%%%%%%%%%%%%%%		EJERCICIOS				%%%%%%%%%%%%%%%%%%%%%
%%%%%%%%%%%%%%%%%%%%%%%%%%%%%%%%%%%%%%%%%%%%%%%%%%%%%%%%%%%%%%%%%%%%%%%%%%
\subsection{Explicación e Implementación de Clases}
 
El proyecto se compone de 5 clases Java que interactúan bajo el patrón de diseño orientado a componentes. A continuación se detallan las porciones de código más importantes del desarrollo:

\subsubsection{1. Clase Pregunta}
Almacena la trivia y, fundamentalmente, la explicación pedagógica que permite la retroalimentación.
\begin{lstlisting}[language=Java, caption={Definición de atributos y constructor de Pregunta}]
public class Pregunta {
    private String texto;
    private String[] opciones;
    private int respuestaCorrecta;
    private String explicacion; // Retroalimentacion educativa

    public Pregunta(String texto, String opcion1, String opcion2, String opcion3,
                    int respuestaCorrecta, String explicacion) {
        this.texto = texto;
        this.opciones = new String[]{opcion1, opcion2, opcion3};
        this.respuestaCorrecta = respuestaCorrecta - 1; 
        this.explicacion = explicacion;
    }
    public String getExplicacion() { return explicacion; }
}
\end{lstlisting}

\subsubsection{2. Clase Jugador}
Soporta el almacenamiento de puntos, tiempo acumulado de respuesta y el avatar del jugador:
\begin{lstlisting}[language=Java, caption={Atributos y constructor de Jugador}]
public class Jugador {
    private String nombre;
    private char simbolo;
    private int puntos;
    private Color color;
    private long tiempoTotal;
    private ImageIcon avatarIcono; // Soporte para foto/avatar a color

    public Jugador(String nombre, char simbolo, ImageIcon avatarIcono) {
        this.nombre = nombre;
        this.simbolo = simbolo;
        this.puntos = 0;
        this.tiempoTotal = 0;
        this.avatarIcono = avatarIcono;
        
        if (simbolo == 'X') {
            this.color = new Color(30, 130, 255); // Azul electrico
        } else {
            this.color = new Color(230, 60, 60);  // Rojo intenso
        }
    }
}
\end{lstlisting}

\subsubsection{3. Renderizado de Avatares a Color con Java 2D}
Para solucionar el aspecto monocromático y plano de los emojis por defecto en Windows/Swing, se implementó un motor de dibujo de avatares vectoriales. A continuación se presenta el código que realiza este proceso:

\begin{lstlisting}[language=Java, caption={Renderizado vectorial a color}, numbers=left]
private ImageIcon renderAvatarColorido(String avatarId, int size, Color colorFondo) {
    BufferedImage img = new BufferedImage(size, size, BufferedImage.TYPE_INT_ARGB);
    Graphics2D g2d = img.createGraphics();
    g2d.setRenderingHint(RenderingHints.KEY_ANTIALIASING, RenderingHints.VALUE_ANTIALIAS_ON);

    // Fondo circular con degradado suave
    GradientPaint gp = new GradientPaint(0, 0, colorFondo.brighter(), size, size, colorFondo.darker());
    g2d.setPaint(gp);
    g2d.fillOval(2, 2, size - 5, size - 5);

    // Borde circular blanco elegante
    g2d.setColor(Color.WHITE);
    g2d.setStroke(new BasicStroke(2.0f));
    g2d.drawOval(2, 2, size - 5, size - 5);

    int c = size / 2; // Centro

    if (avatarId.equals("🤖")) {
        // Cabeza metalica
        g2d.setColor(new Color(175, 189, 205));
        g2d.fillRoundRect(c - 16, c - 10, 32, 24, 6, 6);
        // Antena dorada
        g2d.setColor(new Color(255, 195, 0));
        g2d.setStroke(new BasicStroke(2.5f));
        g2d.drawLine(c, c - 10, c, c - 18);
        g2d.fillOval(c - 3, c - 22, 6, 6);
        // Ojos LED cyan
        g2d.setColor(new Color(0, 255, 255));
        g2d.fillOval(c - 10, c - 4, 6, 6);
        g2d.fillOval(c + 4, c - 4, 6, 6);
        // Boca
        g2d.setColor(new Color(75, 80, 95));
        g2d.drawRect(c - 8, c + 6, 16, 4);
    } else if (avatarId.equals("🐱")) {
        // Gato anaranjado
        g2d.setColor(new Color(255, 165, 55));
        // Orejas y rostro...
        g2d.fillOval(c - 16, c - 8, 32, 23);
    }
    // (Demas avatares vectoriales se omiten por espacio)
    g2d.dispose();
    return new ImageIcon(img);
}
\end{lstlisting}

\paragraph{Explicación del motor gráfico}
\begin{itemize}
	\item \textbf{Líneas 2--4:} Se crea un \texttt{BufferedImage} transparente y se activa el suavizado (\texttt{ANTIALIASING}) para evitar bordes pixelados.
	\item \textbf{Líneas 6--8:} Se define un gradiente dinámico a partir del color del jugador, rellenando el fondo con un círculo.
	\item \textbf{Líneas 10--12:} Se traza el borde circular blanco que brinda un aspecto elegante de medalla.
	\item \textbf{Líneas 16--29:} Se dibujan primitivas del robot (cabeza, antena, ojos LED) utilizando colores vibrantes e independientes del sistema operativo.
\end{itemize}

\begin{figure}[H]
	\centering
	\includegraphics[width=0.85\textwidth,keepaspectratio]{img/diagrama_de_flujo.png}
	\caption{Flujo de interacción y registro de avatares vectoriales en el juego}
\end{figure}

\subsubsection{4. Lógica de Desempate y Resultados}
Al culminar el juego (tablero lleno), si hay empate en puntos, se aplica el criterio de menor tiempo acumulado de respuesta (velocidad mental):

\begin{lstlisting}[language=Java, caption={Lógica de decisión del ganador y desempate}]
if (p1 > p2) {
    ganadorNombre = jugador1.getNombre();
    motivoGanador = "tiene mas puntos (" + p1 + " vs " + p2 + ")";
} else if (p2 > p1) {
    ganadorNombre = jugador2.getNombre();
    motivoGanador = "tiene mas puntos (" + p2 + " vs " + p1 + ")";
} else if (t1 < t2) {
    ganadorNombre = jugador1.getNombre();
    motivoGanador = "respondio mas rapido ("
        + String.format("%.1f", t1 / 1000.0) + "s vs "
        + String.format("%.1f", t2 / 1000.0) + "s)";
}
\end{lstlisting}

\clearpage
 
\subsection{Compilación y Ejecución}
La compilación se realiza directamente con la herramienta \texttt{javac} nativa y la ejecución con la máquina virtual \texttt{java}:

\begin{lstlisting}[language=bash, caption={Comandos para compilar y ejecutar el juego}]
$ javac -encoding UTF-8 *.java
$ java Main
\end{lstlisting}

\begin{figure}[H]
	\centering
	\includegraphics[width=0.85\textwidth,keepaspectratio]{img/codigo.png}
	\caption{Evidencia de compilación exitosa y ejecución de Main en Windows}
\end{figure}

\clearpage
 
%%%%%%%%%%%%%%%%%%%%%%%%%%%%%%%%%%%%%%%%%%%%%%%%%%%%%%%%%%%%%%%%%%%%%%%%%%
%%%%%%%%%%%%%%%%%%%%%			TAREA				%%%%%%%%%%%%%%%%%%%%%
%%%%%%%%%%%%%%%%%%%%%%%%%%%%%%%%%%%%%%%%%%%%%%%%%%%%%%%%%%%%%%%%%%%%%%%%%%
\section{Tarea}
 
\subsection{Pregunta y Respuesta}
\begin{itemize}
	\item \textbf{¿Cómo podríamos saber con exactitud cuál de todas las propuestas algorítmicas es la más óptima? Explique.}
	
	\textit{Respuesta:} Para determinar la optimización con exactitud, se deben evaluar dos factores fundamentales mediante el análisis de algoritmos:
	
	1. \textbf{Complejidad Temporal ($O$):} Se evalúa mediante la notación Big O el número de operaciones que realiza el algoritmo con relación al tamaño de la entrada ($n$). Una búsqueda en un \texttt{HashMap} es $O(1)$ promedio, siendo más óptima que una búsqueda secuencial en una lista que es $O(n)$.
	
	2. \textbf{Complejidad Espacial:} Evalúa el consumo de memoria del algoritmo. Por ejemplo, en Swing y manejo de gráficos, dibujar elementos vectoriales al vuelo en memoria (\texttt{Graphics2D}) evita la carga de recursos de disco, optimizando drásticamente el espacio y evitando I/O bloqueante.
\end{itemize}
 
\clearpage
 
%%%%%%%%%%%%%%%%%%%%%%%%%%%%%%%%%%%%%%%%%%%%%%%%%%%%%%%%%%%%%%%%%%%%%%%%%%
%%%%%%%%%%%%%%%%%%%%%		ENTREGABLES			%%%%%%%%%%%%%%%%%%%%%
%%%%%%%%%%%%%%%%%%%%%%%%%%%%%%%%%%%%%%%%%%%%%%%%%%%%%%%%%%%%%%%%%%%%%%%%%%
\subsection{Entregables}
El proyecto se encuentra alojado en el repositorio GitHub correspondiente a la asignatura, con la siguiente estructura de archivos:
 
\begin{lstlisting}[style=ascii-tree]
TresEnRayaGUI/
|--- .gitignore
|--- README.md
|--- Main.java
|--- VentanaJuego.java
|--- Jugador.java
|--- Pregunta.java
|--- GestorPreguntas.java
|--- latex/
|       |--- informe.tex
|       |--- informe.pdf
|       |--- img/
|               |--- logo_internacional_secundario_color.png
|               |--- IMAGEN1.png
|               |--- diagrama_de_flujo.png
|               |--- codigo.png
\end{lstlisting}
 
%%%%%%%%%%%%%%%%%%%%%%%%%%%%%%%%%%%%%%%%%%%%%%%%%%%%%%%%%%%%%%%%%%%%%%%%%%
%%%%%%%%%%%%%%%%%%%%%			RÚBRICAS			%%%%%%%%%%%%%%%%%%%%%
%%%%%%%%%%%%%%%%%%%%%%%%%%%%%%%%%%%%%%%%%%%%%%%%%%%%%%%%%%%%%%%%%%%%%%%%%%
\section{Rúbricas}
 
\subsection{Evaluación del Proyecto}
\begin{table}[H]
	\caption{Rúbrica de evaluación - Contenido del Proyecto}
	\setlength{\tabcolsep}{0.5em}
	{\renewcommand{\arraystretch}{1.5}
	\begin{tabular}{|p{2.7cm}|p{7cm}|x{1.3cm}|p{1.2cm}|p{1.5cm}|p{1.1cm}|}
		\hline
		\multicolumn{2}{|c|}{Contenido y demostración} & Puntos & Checklist & Estudiante & Profesor \\
		\hline
		\textbf{1. Interfaz Swing} & Tablero funcional 3x3 a color con botones adaptados y barra de estado de turnos. & 4 & SI & 4 & \\
		\hline
		\textbf{2. Banco Trivia} & Banco de preguntas sobre POO y Swing con explicaciones de retroalimentación. & 4 & SI & 4 & \\
		\hline
		\textbf{3. Carga Avatares} & Diálogo de registro para fotos locales e iconos vectoriales dinámicos de color. & 4 & SI & 4 & \\
		\hline
		\textbf{4. Sonidos} & Sonido MIDI nativo para acierto, fallo y victoria integrado. & 2 & SI & 2 & \\
		\hline
		\textbf{5. Desempate} & Algoritmo de desempate por milisegundos en caso de empate en tablero. & 2 & SI  & 2 & \\
		\hline
		\textcolor{red}{\textbf{6. Madurez}} & \textcolor{blue}{El Informe muestra una evolución de la madurez del código con explicaciones precisas del motor de renderizado y el temporizador.} & \textcolor{blue}{4} & SI & 4 & \\
		\hline
		\multicolumn{2}{|c|}{\textbf{Total}} & 20 && 20 & \\
		\hline
	\end{tabular}
	}
\end{table}
 
\clearpage
 
%%%%%%%%%%%%%%%%%%%%%%%%%%%%%%%%%%%%%%%%%%%%%%%%%%%%%%%%%%%%%%%%%%%%%%%%%%
%%%%%%%%%%%%%%%%%%%%%		REFERENCIAS			%%%%%%%%%%%%%%%%%%%%%
%%%%%%%%%%%%%%%%%%%%%%%%%%%%%%%%%%%%%%%%%%%%%%%%%%%%%%%%%%%%%%%%%%%%%%%%%%
\section{Referencias}
\begin{itemize}
	\item Documentación oficial de Java SE: \url{https://docs.oracle.com/javase/8/docs/api/}
	\item Tutorial de Java Swing y Graphics2D: \url{https://docs.oracle.com/javase/tutorial/uiswing/}
	\item Especificación javax.sound.midi: \url{https://docs.oracle.com/javase/8/docs/technotes/guides/sound/}
\end{itemize}
 
\printindex
 
\end{document}
